% page 599 du fac-simile (image p-598 du scan)
% bloc 180.3 x 242.0 mm ; marge gauche 21.7 mm
\pgc{0.000mm}{0.000mm}{
\l{\cel{60}591}
\l{}
\l{\cel{6}\sou{transpirar}. (\sou{netrans}.) I. Pasar tra la pori di korpo, forme di}
\l{\cel{9}guteti qui kelke aspektas quale sudoro. - II. (\sou{metaf.})}
\l{\cel{9}(\sou{pri informi}). Diskonocesar pro la nediskreteso di ula personi}
\l{\cel{9}o di ula persono. - DEFIS.}
\l{}
\l{\cel{6}\sou{transportar}. (\sou{trans.}) Portar de ca a ta loko. - DEFIRS.}
\l{}
\l{\cel{6}\sou{transubstanciar}. (\sou{teol. katol.})\cel{2}Agar la transmuto di la pano}
\l{\cel{9}e di la vino a la korpo e la sango di Iesu Kristo.- DEFIS.}
\l{}
\l{\cel{6}\sou{transversa}. (\sou{pri korpo)}.\cel{2}Qua, de un ek sua extremaji a la altra,}
\l{\cel{9}trairas lineo, direciono, krucume. - DEFIS.}
\l{}
\l{\cel{6}\sou{trapo}. Pordo horizontala qua klozas aperturo facita en la plan-}
\l{\cel{9}karo e quan onu levas od abasas segun-vole. - EFIS.}
\l{}
\l{\cel{6}\sou{trapezo}. I. (\sou{geom.)}\cel{2}Quadrilatero di qua du lateri esas inter-}
\l{\cel{9}paralela, ma ne-interegala. - II. Implemento di gimnastiko}
\l{\cel{9}ek bastono qua povas portar la pezo di persono, suspendita}
\l{\cel{9}per du kordi ringoza a la portiko. - DEFIRS.}
\l{}
\l{\cel{6}\sou{trapezoedro}. I. Solido di qua la edri esas trapezoida. - II.}
\l{\cel{9}Solido ek duadek-e-quar edri quadrilatera intersimetre. -}
\l{\cel{9}DEFIS.}
\l{}
\l{\cel{6}\sou{trapezoido}. I. (\sou{geom.}) Quadrilatero plana, di qua omna lateri}
\l{\cel{9}esas inter-obliqua. - II. (\sou{anat.}) Muskulo sub-pela qua kovras}
\l{\cel{9}preske tote la regiono dorsala. - DEFIS.}
\l{}
\l{\cel{6}\sou{trasar.}\cel{2}(\sou{trans.)}\cel{3}Indikar per streko, per lineo, la direciono,l}
\l{\cel{9}formo. - F.}
\l{}
\l{\cel{6}\sou{tratar}. (\sou{trans.,ulu}). (\sou{aludante komercisto)}. Emisar mandato di}
\l{\cel{9}qua la pekunio-quanto stipulita esas pagenda da (la tratito)}
\l{\cel{9}po la vari furnisita a lu. - DFIRS.}
\l{}
\l{\cel{6}\sou{traumato}. (\sou{patol.}) Ensemblo de la desordini fiziologiala quin}
\l{\cel{9}produktis ago od efiko violentoza. - DEFIRS.}
\l{}
\l{\cel{6}\sou{traurar}. (\sou{netrans.)}\cel{1}Afliktesar da la morto di persono kara. - DR.}
\l{}
\l{\cel{6}\sou{travo}. (\sou{tekn.)}\cel{1}Mashino en qua onu ligas la kavali, la bovi, e c.,}
\l{\cel{9}sive por hufizar li kande li rezistas, sive por submisar li}
\l{\cel{9}ad ula operaci. - EFIS.}
\l{}
\l{\cel{6}\sou{traveo.}\cel{1}(\sou{arkitakt.)}\cel{2}La spaco inter du suportili (trabi, koloni,}
\l{\cel{9}pilastri, e c.) - F.}
\l{}
\l{\cel{6}\sou{traverso}. (\sou{tekn.)}\cel{2}Peco ek ligno o metalo perpendikla ye la ele-}
\l{\cel{9}menti precipua di konstrukturo, destinita a mantenar la inter-}
\l{\cel{9}disto di ta elementi. - EFIS.}
\l{}
\l{\cel{6}travertino. (\sou{geol.})\cel{2}Tofo kalkoza, harda, griza, uzata por la}
\l{\cel{9}konstrukto di domi, e c. -}
}
